\section{Conclusion}
\label{sec:conclusion}

In this paper, we systematically evaluate ChatGPT(gpt3.5-turbo-0301) for the task of product information extraction. We experiment with different zero-shot prompt designs, which show that a proper prompt design is sufficient to achieve similar performance to that of a fine-tuned NER baseline.
%Todo Alex: Change conclusion to not focus on specific performance improvement, but to state that a similar performance is reached with much less training examples.
By adding demonstrations to the prompts, the results are further improved.
%we are able to further improve F1 score of 94.51\%.
While state-of-the-art models like AVEQA outperform ChatGPT's attribute/value pair extractions, they require far more training examples (71,266 vs. 6) and computation ($\sim$10 hrs.) for fine-tuning.
Our experiments also show that ChatGPT can normalize the extracted attribute/value pairs. These normalizations improve the quality of the extracted attribute/value pairs.
